\documentclass[10pt,a4paper]{amsart}
\usepackage[margin=19mm]{geometry}
\usepackage{amsmath,amssymb,mathtools}
\usepackage[scr=rsfs,cal=euler]{mathalfa}
\usepackage{xfrac}
\usepackage[unicode,hidelinks,bookmarks=false]{hyperref}
\newcommand{\ie}{i.e.\ }
\newcommand{\cf}{cf.\ }
\newcommand{\resp}{respectively\ }
\newcommand{\Subsection}{\subsection}
\providecommand{\nobreakdash}{\nobreak\hbox{-}}
\setlength{\emergencystretch}{2em}
\multlinegap0pt
\usepackage{mathtools}
\usepackage{latexsym}
\newtheorem{theorem}{Theorem}[section]
\newtheorem{algorithm}[theorem]{Algorithm}
\newtheorem{axiom}[theorem]{Axiom}
\newtheorem{case}[theorem]{Case}
\newtheorem{claim}[theorem]{Claim}
\newtheorem{conclusion}[theorem]{Conclusion}
\newtheorem{condition}[theorem]{Condition}
\newtheorem{conjecture}[theorem]{Conjecture}
\newtheorem{corollary}[theorem]{Corollary}
\newtheorem{criterion}[theorem]{Criterion}
\newtheorem{definition}[theorem]{Definition}
\newtheorem{exercise}[theorem]{Exercise}
\newtheorem{lemma}[theorem]{Lemma}
\newtheorem{notation}[theorem]{Notation}
\newtheorem{problem}[theorem]{Problem}
\newtheorem{proposition}[theorem]{Proposition}
\newtheorem{remark}[theorem]{Remark}
\newtheorem{solution}[theorem]{Solution}
\newtheorem{summary}[theorem]{Summary}
\newtheorem{example}[theorem]{Example}
\DeclareMathOperator{\ns}{ns}
\DeclareMathOperator{\nc}{nc}
\DeclareMathOperator{\dn}{dn}
\DeclareMathOperator{\ds}{ds}
\DeclareMathOperator{\cs}{cs}
\DeclareMathOperator{\sn}{sn}
\DeclareMathOperator{\cn}{cn}
\DeclareMathOperator{\nd}{nd}
\DeclareMathOperator{\scJ}{sc}
\DeclareMathOperator{\cd}{cd}
\DeclareMathOperator{\dc}{dc}
\DeclareMathOperator{\sd}{sd}
\DeclareMathOperator{\pq}{pq}
\DeclareMathOperator{\qp}{qp}
\DeclareMathOperator{\pn}{pn}
\DeclareMathOperator{\qn}{qn}
\DeclareMathOperator{\np}{np}
\DeclareMathOperator{\qr}{qr}
\DeclareMathOperator{\rqJ}{rq}
\DeclareMathOperator{\p}{p}
\DeclareMathOperator{\q}{q}
\DeclareMathOperator{\s}{s}
\DeclareMathOperator{\cJ}{c}
\DeclareMathOperator{\dJ}{d}
\DeclareMathOperator{\n}{n}
\begin{document}
\title[Generalized Addition Theorems for the Jacobi Elliptic Functi]{Generalized Addition Theorems for the Jacobi Elliptic Functions}
\author{Efe Gürel}
\date{}
\maketitle
\noindent\emph{Source and licence.} Generalized Addition Theorems for the Jacobi Elliptic Functions. Version arXiv:2608.12399v1 (CC BY 4.0). This material is distributed under the Creative Commons Attribution 4.0 International licence, http://creativecommons.org/licenses/by/4.0/, as recorded in the arXiv OAI licence metadata for this exact version and shown on the abstract page https://arxiv.org/abs/2608.12399v1. Full rights notice: licenses/arxiv-2608.12399-CC-BY-4.0.txt.\par\smallskip
\noindent\textit{Editorial scope.} Counted unit: the original \texttt{theorem} environment \texttt{generaldetthm} at source lines 304--314 together with its complete proof at lines 315--355; the delivered document keeps source lines 227--356 so that every referenced label and every citation resolves inside it. Native editable source is the paper's own concatenated TeX; the AMS wrapper is editorial formatting only. No publisher class, upstream script or unrelated artwork is redistributed.
\section{General Determinant Theorems}
\subsection{Notions and the Main Theorem}

We shall first present the fundamental properties of Jacobi elliptic functions. It is well-known that each one of the Jacobi functions $\pq$ satisfies an algebraic differential equation, the most famous one being $\left(\sn'\right)^2=\left(1-\sn^2\right)\left(1-k^2\sn^2\right)$ as previously stated. We put all these equations together in the form
\begin{align*}
    \left(\pq'\right)^2=A_{\pq}\pq^4+B_{\pq}\pq^2+C_{\pq}
\end{align*}
for constants $A_{\pq},B_{\pq},C_{\pq}$. Explicit forms of these equations are not needed in our investigations and their mere existence will suffice. The reader may consult classic treatises such as \cite{Whittakter} or \cite{dlmf} for these equations. The half-period translations induce a sign change as follows,
\begin{align*}
    \sn(u+2mK+2niK')&=(-1)^{m}\sn u,\\
    \cn(u+2mK+2niK')&=(-1)^{m+n}\cn u,\\
    \dn(u+2mK+2niK')&=(-1)^{n}\dn u.
\end{align*}
Furthermore, the quarter-period translations transform the Jacobi functions into one another. The following table may be found in \cite{dlmf}.
\begin{center}
\begin{table}[h]
    \begin{tabular}{|c|c|c|c|}
    \hline
    Quarter-period $\omega$ & $K$ & $iK'$ & $K+iK'$\\ \hline
    $\sn(u+\omega)$ & $\cd u$ & $\frac{1}{k}\ns u$ & $\frac{1}{k}\dc u$ \\ \hline
    $\cn(u+\omega)$ & $-k' \sd u$ & $-i\frac{1}{k}\ds u$ & -i$\frac{k'}{k}\nc u$ \\ \hline      
    $\dn(u+\omega)$ & $k'\nd  u$ & $-i\cs u$ & $i k'\scJ u$ \\ \hline      
    \end{tabular}\medskip
    
    \caption{Quarter-period Transformations of Jacobi Functions}
\end{table}
\end{center}
These transformations can be used to transform addition theorems of one function into another. We also give the following table of periods and poles with further information in the case where periodicity is expanded to modulo $(4K,4iK')$.

\begin{center}
\begin{table}[h]
\begin{tabular}{|c|c|c|c|}
\hline
Function & Minimal Periods & Poles modulo Periods & Poles modulo $(4K,4iK')$ \\ \hline
$\sn$ & $4K,2iK'$ & $iK',iK'+2K$ & $iK',iK'+2K,3iK',3iK'+2K$ \\ \hline
$\cn$ & $4K,2K+2iK'$ &  $iK',iK'+2K$ & $iK',iK'+2K,3iK',3iK'+2K$ \\ \hline
$\dn$ & $2K,4iK'$ &  $iK',3iK'$ & $iK',iK'+2K,3iK',3iK'+2K$ \\ \hline
$\ns$ & $4K,2iK'$ & $0,2K$ & $0,2K,2iK',2K+2iK'$ \\ \hline
$\cs$ & $2K,4iK'$ & $0,2iK'$ & $0,2K,2iK',2K+2iK'$ \\ \hline
$\ds$ & $4K,2K+2iK'$ & $0,2K$ & $0,2K,2iK',2K+2iK'$ \\ \hline
$\nc$ & $4K,2K+2iK'$ & $K,3K$ & $K,3K,K+2iK',3k+2iK'$ \\ \hline
$\scJ$ & $2K,4iK'$ & $K,K+2iK'$ & $K,3K,K+2iK',3k+2iK'$ \\ \hline
$\dc$ & $4K,2iK'$ & $K,3K$ & $K,3K,K+2iK',3k+2iK'$ \\ \hline
$\nd$ & $2K,4iK'$ & $K+iK',K+3iK'$ & $K+iK',K+3iK',3K+iK',3K+3iK'$ \\ \hline
$\sd$ & $4K,2K+2iK'$ & $K+iK',3K+iK'$ & $K+iK',K+3iK',3K+iK',3K+3iK'$ \\ \hline
$\cd$ & $4K,2iK'$ & $K+iK',3K+iK'$ & $K+iK',K+3iK',3K+iK',3K+3iK'$ \\ \hline
\end{tabular}\medskip
\caption{Jacobi Functions, Their Periods and Poles}
\label{periodtable}
\end{table}
\end{center}
Notice that the sum of the poles modulo minimal periods is always a half-period. More importantly, it is always a period when expanded to modulo $(4K,4iK')$. Let $(\pq)^{(-1)}=1$ be the constant function for convenience. This allows us to generalize so that the functions $\pq^{\alpha}$ and $\pq^{(\beta)}$ are elliptic functions of orders $2\alpha$ and $2\beta+2$ with respect to their minimal periods. However, they are of order $4\alpha$ and $4\beta+4$ with respect to the periods $(4K,4iK')$. We can also divide the functions into period and pole classes. Naturally, a period class is the set of functions having equal minimal periods and can be found as $P_{\p}=P_{\qr}=\{\pn,\np,\qr,\rqJ\}$ for a permutation $\p\q\text{r}$ of $\s\cJ\dJ$. A pole class $Z_{\p}$ is determined as the set of functions that share the common poles when the periodicity is expanded to modulo $(4K,4iK')$. It can be seen that $Z_{\p}=\{\qp, \q\neq \p\}$ for any letter $\p$, as shown in the Table \ref{periodtable}. \medskip
 

Take $r$ of the $12$ Jacobi functions $(\pq)_1,\ldots, (\pq)_r$. Let us have distinct positive integers $\alpha_{\ell,j},\alpha_{\ell,j'}'$ where $\ell=1,\ldots, r$ and $j=1,\ldots, n_\ell, j'=1,\ldots,n'_\ell$. Furthermore, assume $\beta_{\ell,j},\beta_{\ell,j'}'\in \mathbb{Z}^+\cup\{-1\}$ where $\ell=1,\ldots, r$ and $j=1,\ldots, m_\ell, j'=1,\ldots,m'_\ell$. We can order the numbers such that $\alpha_{\ell,1}< \alpha_{\ell,2}< \ldots < \alpha_{\ell,n_\ell}$ and the respective inequalities hold for $\alpha_{\ell,j'}',\beta_{\ell,j},\beta_{\ell,j'}'$. Assume that these variables satisfy the condition
\begin{align}\label{VarCond}
    \sum_{\p\in\{\s,\cJ,\dJ,\n \}} \max_{\ell\in Z_{\p}}\left(4\alpha_ {\ell,n_\ell},4\alpha'_{\ell,n_\ell'},4\beta_{\ell,m_\ell}+4,4\beta_{\ell,m_\ell'}'+4  \right)=1+\sum_{\ell=1}^r n_\ell+m_\ell.
\end{align}
Let us be given \textit{generic} complex numbers $\mu_{\ell,j_1},\nu_{\ell,j_2}$ with $\ell=1,\ldots,r$ and $j_1=1,\ldots, n_\ell', j_2=1,\ldots,m_\ell'$. Finally, take $N=\sum_{\ell=1}^r n_\ell+m_\ell$ \textit{generic} complex numbers $z_1,\ldots, z_N$. The word generic here will be used in an implicit manner to mean parameters not satisfying any special equations. For example, in our investigations we shall make use of the assumption that $\pq(z_h)$ are distinct. This assumption imposes a very relaxed condition on $z_h$ as the function $\pq$ takes each complex value twice in a fundamental parallelogram. As another example, we shall calculate orders of elliptic functions formed by linear combinations of powers and derivatives of $\pq$. We will then take the maximum order for a pole as granted. This will be true for \textit{generic} values of parameters. Indeed, if the leading terms around a pole, say coming from $(\pq)_\ell^{\alpha}$ and $(\pq)_\ell^{(\beta)}$ or more, cancel each other to yield a non-maximal order, this would impose a linear relation between the parameters. This relation does not hold in general and may be avoided by an appropriate choice of the parameters. However, our results will hold for \textit{all} values of the parameters by the means of analytic continuation or taking limits. In fact, non-generic parameters are associated with degenerate cases of our formula, such as when a pole is encountered or a double root is present. Our proofs can be modified to include these degenerate cases. However, this is unnecessary as the process of analytic continuation provides a much simpler solution to this issue. The reader may consult \cite{Gürel} for examples of these phenomena. \medskip

We have
\begin{align*}
    \begin{vmatrix}
(\pq)_1^{\alpha_{1,1}}(z_1) & \cdots & (\pq)_1^{\alpha_{1,n_1}}(z_1) & (\pq)_1^{(\beta_{1,1})}(z_1) & \cdots & (\pq)_1^{(\beta_{1,m_1})}(z_1) & \cdots &  (\pq)_r^{(\beta_{r,m_r})}(z_1) \\
 \vdots& \cdots & \cdots &  \cdots & \cdots & \cdots & \cdots &  \vdots \\
 (\pq)_1^{\alpha_{1,1}}(z_N)&\cdots  & \cdots & \cdots & \cdots & \cdots & \cdots &  (\pq)_r^{(\beta_{r,m_r})}(z_N)
\end{vmatrix}\neq 0
\end{align*}
for generic values of $z_1,\ldots,z_N$. This guarantees the existence of complex numbers $\lambda_{\ell,j_1},\kappa_{\ell,j_2}$ with $j_1=1,\ldots,n_\ell, j_2=1,\ldots, m_\ell$ such that the equation
\begin{align*}
    &\lambda_{1,1}(\pq)_1^{\alpha_{1,1}}(z_h)+\ldots+\lambda_{1,n_1}(\pq)_1^{\alpha_{1,n_1}}(z_h)+\kappa_{1,1}(\pq)_1^{(\beta_{1,1})}(z_h)+\ldots+\kappa_{1,m_1}(\pq)_1^{(\beta_{1,m_1})}(z_h)+\ldots\\
+&\lambda_{r,1}(\pq)_r^{\alpha_{r,1}}(z_h)+\ldots+\lambda_{r,n_r}(\pq)_r^{\alpha_{r,n_r}}(z_h)+\kappa_{r,1}(\pq)_r^{(\beta_{r,1})}(z_h)+\ldots+\kappa_{r,m_r}(\pq)_r^{(\beta_{r,m_r})}(z_h)\\
=&\mu_{1,1}(\pq)_1^{\alpha_{1,1}'}(z_h)+\ldots+\mu_{1,n_1'}(\pq)_1^{\alpha_{1,n_1'}'}(z_h)+\nu_{1,1}(\pq)_1^{(\beta'_{1,1})}(z_h)+\ldots+\nu_{1,m_1'}(\pq)_1^{(\beta'_{1,m_1'})}(z_h)+\ldots\\
+&\mu_{r,1}(\pq)_r^{\alpha_{r,1}'}(z_h)+\ldots+\mu_{r,n_r'}(\pq)_r^{\alpha_{r,n_r'}'}(z_h)+\nu_{r,1}(\pq)_r^{(\beta'_{r,1})}(z_h)+\ldots+\nu_{r,m_r'}(\pq)_r^{(\beta'_{r,m_r'})}(z_h)
\end{align*}
holds for all $h=1,\ldots,N$. We remark that the constants $\lambda_{\ell,j_1}$ and $\kappa_{\ell,j_2}$ are rational functions of $\pq(z_h),\pq'(z_h)$ for $h=1,\ldots,N$ and the constants $\mu_{\ell,j_1},\nu_{\ell,j_2}$. They can be calculated \textit{explicitly} by Cramer's formula. Finally, let $z$ be defined by $z+z_1+\ldots+z_N=0$. Our main theorem follows.


\begin{theorem}\label{generaldetthm}
    The following equation holds,
    \begin{align*}
        \begin{vmatrix}
\sum_{\ell=1}^r\sum_{j=1}^{n_\ell'} \mu_{\ell,j}(\pq)_\ell^{\alpha_{\ell,j}'}(z_1)+\sum_{\ell=1}^r\sum_{j=1}^{m_\ell'} \nu_{\ell,j}(\pq)_\ell^{(\beta_{\ell,j}')}(z_1)& (\pq)_1^{\alpha_{1,1}}(z_1) &  \cdots & (\pq)_r^{(\beta_{r,m_r})}(z_1)  \\
 \vdots& \vdots & \cdots & \vdots \\
\sum_{\ell=1}^r\sum_{j=1}^{n_\ell'} \mu_{\ell,j}(\pq)_\ell^{\alpha_{\ell,j}'}(z_N)+\sum_{\ell=1}^r\sum_{j=1}^{m_\ell'} \nu_{\ell,j}(\pq)_\ell^{(\beta_{\ell,j}')}(z_N)& (\pq)_1^{\alpha_{1,1}}(z_N) &  \cdots & (\pq)_r^{(\beta_{r,m_r})}(z_N) \\
\sum_{\ell=1}^r\sum_{j=1}^{n_\ell'} \mu_{\ell,j}(\pq)_\ell^{\alpha_{\ell,j}'}(z)+\sum_{\ell=1}^r\sum_{j=1}^{m_\ell'} \nu_{\ell,j}(\pq)_\ell^{(\beta_{\ell,j}')}(z)& (\pq)_1^{\alpha_{1,1}}(z) &  \cdots & (\pq)_r^{(\beta_{r,m_r})}(z)
\end{vmatrix}=0.
    \end{align*}    
\end{theorem}

\begin{proof}
    Let us consider the function 
    \begin{align*}
        \psi(u)=&\lambda_{1,1}(\pq)_1^{\alpha_{1,1}}(u)+\ldots+\lambda_{1,n_1}(\pq)_1^{\alpha_{1,n_1}}(u)+\kappa_{1,1}(\pq)_1^{(\beta_{1,1})}(u)+\ldots+\kappa_{1,m_1}(\pq)_1^{(\beta_{1,m_1})}(u)+\ldots\\
+&\lambda_{r,1}(\pq)_r^{\alpha_{r,1}}(u)+\ldots+\lambda_{r,n_r}(\pq)_r^{\alpha_{r,n_r}}(u)+\kappa_{r,1}(\pq)_r^{(\beta_{r,1})}(u)+\ldots+\kappa_{r,m_r}(\pq)_r^{(\beta_{r,m_r})}(u)\\
-&\mu_{1,1}(\pq)_1^{\alpha_{1,1}'}(u)-\ldots-\mu_{1,n_1'}(\pq)_1^{\alpha_{1,n_1'}'}(u)-\nu_{1,1}(\pq)_1^{(\beta'_{1,1})}(u)-\ldots-\nu_{1,m_1'}(\pq)_1^{(\beta'_{1,m_1'})}(u)-\ldots\\
-&\mu_{r,1}(\pq)_r^{\alpha_{r,1}'}(u)-\ldots-\mu_{r,n_r'}(\pq)_r^{\alpha_{r,n_r'}'}(u)-\nu_{r,1}(\pq)_r^{(\beta'_{r,1})}(u)-\ldots-\nu_{r,m_r'}(\pq)_r^{(\beta'_{r,m_r'})}(u).
    \end{align*}
It is an elliptic function of periods $4K,4iK'$ with the order
\begin{align*}
        \sum_{\p\in\{\s,\cJ,\dJ,\n \}} \max_{\ell\in Z_{\p}}\left(4\alpha_ {\ell,n_\ell},4\alpha'_{\ell,n_\ell'},4\beta_{\ell,m_\ell}+4,4\beta_{\ell,m_\ell'}'+4  \right)=N+1
\end{align*}
for generic parameters. Assuming that the numbers $z_h$ are distinct, we have $N$ distinct zeros already. Let $w$ be the last zero. By Table \ref{periodtable}, the poles sum to zero modulo $(4K,4iK')$. Thus Abel's theorem gives us $w+z_1+z_2+\ldots+z_N\equiv 0 \bmod (4K,4iK')$ and $w\equiv z \bmod (4K,4iK')$. Therefore we have
\begin{align*}
        &\lambda_{1,1}(\pq)_1^{\alpha_{1,1}}(z)+\ldots+\lambda_{1,n_1}(\pq)_1^{\alpha_{1,n_1}}(z)+\kappa_{1,1}(\pq)_1^{(\beta_{1,1})}(z)+\ldots+\kappa_{1,m_1}(\pq)_1^{(\beta_{1,m_1})}(z)+\ldots\\
+&\lambda_{r,1}(\pq)_r^{\alpha_{r,1}}(z)+\ldots+\lambda_{r,n_r}(\pq)_r^{\alpha_{r,n_r}}(z)+\kappa_{r,1}(\pq)_r^{(\beta_{r,1})}(z)+\ldots+\kappa_{r,m_r}(\pq)_r^{(\beta_{r,m_r})}(z)\\
=&\mu_{1,1}(\pq)_1^{\alpha_{1,1}'}(z)+\ldots+\mu_{1,n_1'}(\pq)_1^{\alpha_{1,n_1'}'}(z)+\nu_{1,1}(\pq)_1^{(\beta'_{1,1})}(z)+\ldots+\nu_{1,m_1'}(\pq)_1^{(\beta'_{1,m_1'})}(z)+\ldots\\
+&\mu_{r,1}(\pq)_r^{\alpha_{r,1}'}(z)+\ldots+\mu_{r,n_r'}(\pq)_r^{\alpha_{r,n_r'}'}(z)+\nu_{r,1}(\pq)_r^{(\beta'_{r,1})}(z)+\ldots+\nu_{r,m_r'}(\pq)_r^{(\beta'_{r,m_r'})}(z).
\end{align*}
Writing these equations in the matrix form gives us
\begin{align*}
    &\begin{pmatrix}
\sum_{\ell=1}^r\sum_{j=1}^{n_\ell'} \mu_{\ell,j}(\pq)_\ell^{\alpha_{\ell,j}'}(z_1)+\sum_{\ell=1}^r\sum_{j=1}^{m_\ell'} \nu_{\ell,j}(\pq)_\ell^{(\beta_{\ell,j}')}(z_1)& (\pq)_1^{\alpha_{1,1}}(z_1) &  \cdots & (\pq)_r^{(\beta_{r,m_r})}(z_1)  \\
 \vdots& \vdots & \cdots & \vdots \\
\sum_{\ell=1}^r\sum_{j=1}^{n_\ell'} \mu_{\ell,j}(\pq)_\ell^{\alpha_{\ell,j}'}(z_N)+\sum_{\ell=1}^r\sum_{j=1}^{m_\ell'} \nu_{\ell,j}(\pq)_\ell^{(\beta_{\ell,j}')}(z_N)& (\pq)_1^{\alpha_{1,1}}(z_N) &  \cdots & (\pq)_r^{(\beta_{r,m_r})}(z_N) \\
\sum_{\ell=1}^r\sum_{j=1}^{n_\ell'} \mu_{\ell,j}(\pq)_\ell^{\alpha_{\ell,j}'}(z)+\sum_{\ell=1}^r\sum_{j=1}^{m_\ell'} \nu_{\ell,j}(\pq)_\ell^{(\beta_{\ell,j}')}(z)& (\pq)_1^{\alpha_{1,1}}(z) &  \cdots & (\pq)_r^{(\beta_{r,m_r})}(z)
\end{pmatrix}\\
\cdot&\begin{pmatrix}
-1 \\
 \lambda_{1,1}\\
\vdots\\
\kappa_{r,m_r}
\end{pmatrix}=\begin{pmatrix}
\psi(z_1) \\
\vdots\\
\psi(z_N)\\
\psi(z)
\end{pmatrix}=0.
\end{align*}
Since the vector on the left-hand side is non-zero, the matrix on the left-hand side is singular and has a vanishing determinant. This concludes the proof.
\end{proof}


\begin{thebibliography}{99}
\bibitem[Gürel]{Gürel} Gürel, E. A recipe for obtaining algebraic addition theorems of the Weierstrass elliptic function. \textit{Integral Transforms and Special Functions}, (2026), 1-12.
\bibitem[Whittakter]{Whittakter} Whittaker, E. T., \& Watson, G. N. A Course of Modern Analysis. Courier Dover Publications. (2020).
\bibitem[dlmf]{dlmf} Olver, F. W. (Ed.). NIST Handbook of Mathematical Functions hardback and CD-ROM. Cambridge University Press. (2010).
\end{thebibliography}
\end{document}
